% problem-000012 1 元素周期与氧化还原
\section*{1. 元素周期与氧化还原}
\textit{下列为分问。}

\subsection*{1-1}
离⼦化合物 $\mathbf { A } _ { 2 } \mathbf { B }$ 由四种元素组成，⼀种为氢，另三种为第⼆周期元素。正、负离⼦皆由两种原⼦构成且均呈正四⾯体构型。写出这种化合物的化学式。

\subsection*{1-2}
对碱⾦属Li、Na、K、Rb和Cs，随着原⼦序数增加以下哪种性质的递变不是单调的？简述原因。

(a)熔沸点 (b)原⼦半径 (c)晶体密度 (d)第⼀电离能

\subsection*{1-3}
保险粉 $\mathrm { ( N a _ { 2 } S _ { 2 } O _ { 4 } { \cdot } 2 H _ { 2 } O ) }$ 是重要的化⼯产品，⽤途⼴泛，可⽤来除去废⽔ $\mathrm { ( p H \approx 8 ) }$ 中的Cr(Ⅵ)，所得含硫产物中硫以S(Ⅳ)存在。写出反应的离⼦⽅程式。

\subsection*{1-4}
化学合成的成果常需要⼀定的时间才得以应⽤于⽇常⽣活。例如，化合物A合成于1929年，⾄1969年才被⽤作⽛膏的添加剂和补⽛填充剂成分。A是离⼦晶体，由NaF和 $\mathrm { N a P O _ { 3 } }$ 在熔融状态下反应得到。它易溶于⽔，阴离⼦⽔解产⽣氟离⼦和对⼈体⽆毒的另⼀种离⼦。

\subsection*{1-4}
-1 写出合成A的反应⽅程式。

\subsection*{1-4}
-2 写出A中阴离⼦⽔解反应的离⼦⽅程式。
